\chapter{Proofs for Small-Clause Optimizations}
\label{annex:p12-opt-small-clauses}

\begin{annexitem}{annex:p12.opt-renaming}{thm:p12-opt-renaming}{Interpretation of a new predicate}
\begin{proof}
Fix a \(\Sigma\)-structure \(M\). Set
\[
 R^{M'}=\{\bar d\in D^{|\bar y|}:M,s[\bar y/\bar d]\models A\}.
\]
This is independent of the values of variables outside \(\bar y\):
term evaluation depends only on its occurring variables, by term
induction; atom truth therefore does too; each Boolean operation
preserves this property; and a quantifier removes its bound variable
from the relevant free-variable set by ranging through all its
values. This proves the required free-variable dependence by
formula induction. Thus the displayed set is well defined.
For a nullary predicate \(D^0\) is the singleton empty tuple,
so the definition includes that case.

In this expansion \(R(\bar y)\) and \(A\) agree at all assignments.
The defining sentence holds in every polarity case, and contextual
equivalence proves the forward assertion.
Conversely, a model of the renamed sentence has
\(R(\bar y)\to A\) at all assignments in the positive case, or
\(A\to R(\bar y)\) in the negative case, or equality of truth
values in the zero case. Contextual replacement therefore takes
truth of the renamed context to truth of \(F\). Its evaluation
uses only \(\Sigma\), so it is true in the reduct.
\end{proof}
\end{annexitem}

\begin{annexitem}{annex:p12.opt-counts}{thm:p12-opt-counts}{Every count recurrence}
\begin{proof}
An atom and its negation each produce one clause. Negation
interchanges the positive and negative translations.
Conjunction concatenates two clause lists and hence adds their
lengths. Disjunction distributes each clause of the first list
against each clause of the second, forming their Cartesian product;
its length is the product. De Morgan's laws therefore give,
respectively, \(q_1q_2\) and \(q_1+q_2\) on the negative side.
Implication translates as \(\neg A_1\lor A_2\); its negation as
\(A_1\land\neg A_2\). These give \(q_1p_2\) and \(p_1+q_2\).
For a positive biconditional use
\((\neg A_1\lor A_2)\land(A_1\lor\neg A_2)\), obtaining
\(q_1p_2+p_1q_2\). For a negative one use
\((A_1\lor A_2)\land(\neg A_1\lor\neg A_2)\), obtaining
\(p_1p_2+q_1q_2\). Eliminating an existential quantifier replaces
terms, and retaining a universal quantifier changes the prefix;
neither changes a list length. Quantifier duality gives the same
conclusion for \(q\). Structural induction proves the table.

For the numerical bound, an atom satisfies \(1\le2^N\).
A unary node inherits a child count. At a binary node put
\(N=1+N_1+N_2\). Each product is bounded by
\(2^{N_1+N_2}\), and a sum of two such products by \(2^N\).
A sum of child counts is also at most \(2^N\), since
\(2^{N_1}+2^{N_2}\le2^{N_1+N_2+1}\).
These are all the table operations. Storing both values once at
each node avoids repeated recursion. Addition of \(O(N)\)-bit
integers costs \(O(N)\) bit operations and elementary multiplication
costs \(O(N^2)\); over \(N\) nodes, \(O(N^3)\) is a sufficient
polynomial upper bound. This bound concerns computing the count,
not printing the potentially exponentially many clauses it counts.
\end{proof}
\end{annexitem}

\begin{annexitem}{annex:p12.opt-coefficients}{thm:p12-opt-coefficients}{Coefficients and the saving inequality}
\begin{proof}
At the root take \(a=1,b=c=0\).
Suppose the affine identity holds at a parent.
At a conjunction with selected first child \(B\), its contribution is
\[
 a(p(B)+p_2)+bq(B)q_2+c
   =a p(B)+(bq_2)q(B)+(ap_2+c).
\]
At a disjunction it is
\(ap(B)p_2+b(q(B)+q_2)+c\), giving coefficients
\((ap_2,b)\) and constant \(bq_2+c\).
At an implication's first child it is
\(aq(B)p_2+b(p(B)+q_2)+c\), giving
\((b,ap_2)\); at its second child it is
\(aq_1p(B)+b(p_1+q(B))+c\), giving \((aq_1,b)\).
At a biconditional its contribution is
\[
 a[p(B)q_2+q(B)p_2]+b[p(B)p_2+q(B)q_2]+c,
\]
giving \((aq_2+bp_2,ap_2+bq_2)\).
The other child in each symmetric connective gives the same formulas
with indices interchanged. Negation exchanges the coefficients;
quantifiers preserve them. Every new constant is nonnegative.
Induction proves the affine formula.

All counts are at least one. Starting at \((1,0)\), a nonzero
coefficient remains positive through multiplication by a sibling
count. Negation or an antecedent exchanges the zero/nonzero roles;
a biconditional makes both coefficients positive. This proves
the polarity assertions.
Replacing the occurrence by an atom changes its affine contribution
from \(ap+bq+c\) to \(a+b+c\).
The additional definition costs \(p\), \(q\), or \(p+q\),
respectively in the positive, negative or zero case. Therefore the
three nonincrease conditions are \(ap\ge a+p\),
\(bq\ge b+q\), and \(ap+bq\ge a+b+p+q\).
Subtract and factor to obtain the displayed inequalities.
\end{proof}
\end{annexitem}

\begin{annexitem}{annex:p12.opt-capped}{thm:p12-opt-capped}{Capped arithmetic and an explicit finite construction}
\begin{proof}
For nonnegative integers \(u,v\), direct cases give
\[
 \min(u+v,3)=\min(\min(u,3)+\min(v,3),3),
\]
\[
 \min(uv,3)=\min(\min(u,3)\min(v,3),3).
\]
For the product, a zero factor gives zero on both sides. Otherwise
both factors are positive; a factor at least three makes both
capped products at least three, and factors below three are
represented exactly. The sum identity follows because either
summand at least three saturates both sums, or both are exact.
Induction in the two tables proves that every stored value is the
correct capped count or coefficient.

In a single-sign test a product of nonnegative integers is at least
one exactly when both are at least one. For the zero-sign test the
two terms are nonnegative. Whether \((a-1)(p-1)\) is zero, one,
or at least two is determined exactly by the capped positive
inputs \(a,p\): a value one gives zero; two values two give one;
otherwise two inputs at least two with one at least three give
at least two. Thus capping these products at two preserves whether
their sum reaches two. The negative-sign case is identical.
The explicit threshold table follows by separating the cases that
one input is one from the case both are at least two.

Each node requires boundedly many operations and stores boundedly
many values, proving the linear bound for one pass.
A polynomial implementation uses the markers in
Definition~\ref{def:p12-opt-schedule}. A failed test removes one marker.
An accepted replacement also removes exactly the visited marker:
its descendant markers move to the single copied subformula in the
new definition. They are not present in the replaced occurrence,
which is now a new atom. The new surrounding syntax has no markers,
and binder renaming preserves the remaining markers. Consequently
the number of unvisited markers decreases by one at every visit,
so there are at most \(N\) visits.

At an accepted step, the freshly recomputed coefficients are those
of the current whole formula, including every earlier definition.
The exact saving theorem therefore proves nonincrease for that
whole formula. Its expansion and reduct maps compose with those
of every previous step, proving satisfiability preservation for
the finite sequence. Failed tests change no formula. Induction
over the visits establishes both claimed invariants.

Store each defining formula once, rather than recursively
reinserting its meaning or expanding its biconditionals mid-traversal.
There are at most \(N\) copied original subtrees, each with at most
\(N\) nodes and at most \(N\) free-variable arguments per new
name, giving a polynomial intermediate representation. Recomputing
on those representations is consequently polynomial. No claim of
termination is made for an instruction to rename new definitions
indefinitely.

Here is the complete, independent output construction asserted in
the theorem. For each original nonliteral occurrence \(A\), let
\(R_A\) have its free variables as arguments. In each definition
replace a nonliteral child by its name and retain literal children.
For a binary connective the clause conversion of
\(R_A\leftrightarrow(B\circ C)\), where \(B,C\) are names or
literals, has at most four clauses of at most three literals.
Conjunction and disjunction need three; implication also needs
three; biconditional needs four. For a negation, two suffice.
A quantifier definition has, with the argument lists suppressed,
the form \(R\leftrightarrow QxS\).
For \(Q=\forall\) its clauses are
\[
 \neg R(Z)\lor S(Z,x),\qquad
 R(Z)\lor\neg S(Z,f(Z));
\]
for \(Q=\exists\) they are
\[
 \neg R(Z)\lor S(Z,f(Z)),\qquad
 R(Z)\lor\neg S(Z,x).
\]
Here \(S\) can have a sublist of the displayed arguments, and
\(f\) is fresh for this definition. Any vacuous binder is simply
removed. Each quantified child can be a literal, with the same
substitution applied inside its terms. Names and witnesses for
different occurrences are distinct. There are at most \(4N+1\)
clauses, including the root assertion.

Before introducing the witnesses, interpret every name as the
relation defined by its original subformula. All shallow
equivalences then hold simultaneously, by induction from children
to parents. Conversely the equivalences determine these meanings
by the same induction; the root assertion implies the original
sentence. For the displayed quantified clauses, a model of
\(R\leftrightarrow\forall xS\) chooses a counterexample when \(R\)
is false and an arbitrary value when it is true; a model of
\(R\leftrightarrow\exists xS\) chooses a witness when \(R\) is
true and an arbitrary value when false. This gives the required
total \(f\) for every tuple \(Z\). The two clauses imply the
respective equivalence in the reduct by the same witness and
counterexample reasoning. Finitely many distinct new symbols can
be interpreted simultaneously.

Each clause has at most a constant number of literal occurrences.
A literal before substitution has at most \(O(N)\) symbols:
it is either a name with at most \(N\) arguments or an original
literal of size at most \(N\). A displayed witness term has at
most \(N+1\) symbols and no previously introduced witness occurs
inside it, because each shallow quantifier definition is processed
separately. Replacing at most \(N\) occurrences therefore gives
\(O(N^2)\) symbols per clause and \(O(N^3)\) overall.
This deliberately explicit bound does not identify this construction
with every historical renaming strategy.
\end{proof}
\end{annexitem}

\begin{annexitem}{annex:p12.opt-sharing}{thm:p12-opt-sharing}{Instances of a common name}
\begin{proof}
Interpret \(R(\bar d)\) by the truth of \(A\) at \(\bar y=\bar d\),
as in the single-name proof. The substitution lemma makes every
\(R(\bar y)\sigma_i\) equivalent to \(A\sigma_i\).
That lemma follows by induction: variable and function evaluation
commute with assigning the substituted terms; atoms inherit term
equality; Boolean operations preserve truth equality; and at a
quantifier, prior renaming avoids capture so that assignment of
its fresh bound variable commutes with the substitution.
Thus the expansion satisfies all instances simultaneously.
In a reduct, instantiate the retained definition at the term tuple
\(\bar y\sigma_i\). The resulting directional implication is
exactly the one needed at that occurrence. Replace the finitely
many nonoverlapping occurrences successively using contextual
monotonicity. Other occurrences retain their positions and polarities.
Mixed signs use equivalence and cause no directional conflict.

For syntactic matching, traverse the pattern and target together.
At matching nonvariable symbols check symbol identity and arity,
then recurse on corresponding children. A bound-variable position
must have the same aligned binder identifier. At a substitutable
free-variable occurrence, store the target subtree on its first
visit. At a repeated occurrence, compare the current target subtree
with that stored subtree, rejecting on any difference. Reject a
candidate containing a bound variable that would escape its scope;
under the stipulated capture-safe alignment this is the admissibility
check. If traversal succeeds the collected assignments define the
matching substitution by induction on the pattern; any matching
substitution must pass every comparison, proving the converse.
The target subtrees matched against distinct free-variable
occurrences are disjoint because descent stops at those occurrences.
Comparing a stored subtree to a repeated occurrence costs at most
a constant times the size of that repeated target subtree.
Their sum is at most the whole target size. All other nodes are
visited a bounded number of times. With variable identifiers giving
constant-time access to their recorded matches, the claimed linear
bound follows. Matching modulo logical identities is a different
problem and is not covered by this traversal.
\end{proof}
\end{annexitem}

\begin{annexitem}{annex:p12.opt-standard-skolem}{thm:p12-opt-standard-skolem}{Arbitrary inner and outer occurrences}
\begin{proof}
For an inner step, fix any original structure \(M\).
At every tuple of values of \(\operatorname{FV}(\exists xA)\)
choose a witness for \(A\) if one exists, otherwise an arbitrary
element of \(D\). This defines one total function and makes
\(\exists xA\) and \(A[x/f(\bar y)]\) equivalent at every
assignment. The other free variables of the context cannot affect
the choice because they do not occur free in this subformula.
This is valid even if a variable in \(\bar y\) is bound by an
enclosing existential. Contextual equivalence proves expansion.
The replacement implies the existential in every expansion,
proving the reduct assertion.

For the outer step, first construct a fully outer-Skolemized
version of \(F\), temporarily using one distinct fresh function
for every existential binder. Process them in order with every
ancestor before its descendants. At each step all free variables
of its existential subformula are among the universally bound
ancestors: previously enclosing existentials have already been
replaced by terms in those variables. Choose a witness at each
tuple of all those universal variables whenever one exists,
otherwise a fixed domain element. Extra unused universal
parameters can be ignored by the chosen function.
Exactly the inner witness proof just given, with these padded
arguments, makes the step truth preserving in the selected
expansion. Induction on the finite number of existential binders
therefore constructs a same-domain expansion satisfying the
fully transformed sentence whenever \(M\models F\).

Call the function chosen for the designated occurrence \(g\).
Its arguments are exactly that occurrence's originally enclosing
universal variables. Ancestor existential substitutions may have
changed the body in which its witnesses were chosen, but have
introduced no other free variables into \(g\)'s argument list.
Keep \(g\), and restore all other existential binders in reverse
of the order just used. At a restoration of a binder \(\exists y\),
the current instantiated body implies the restored existential
body, with its former witness term as witness. The \(g\)-terms
remain unchanged because they contain only universal-variable
arguments, never the restored \(y\).
Rectification ensures that no restored binder captures an
argument. For any other previously processed witness term,
reverse order restores descendants before ancestors, so that
its substitution is exactly the one being reversed.
Contextual monotonicity at these positive occurrences preserves
truth at every restoration. The result is precisely \(F\) with
only the designated outer step performed.
Forgetting the other temporary functions gives the asserted
expansion. Conversely \(A[x/g(\bar z)]\) implies \(\exists xA\)
at every assignment in any expansion, so the single-step
sentence implies \(F\) in its reduct.

The distinction about local equivalence is real.
On a two-element domain, \(\exists y\exists x(x\ne y)\)
is true. An outer step at the inner binder uses a constant \(c\).
The whole \(\exists y(c\ne y)\) remains true, but
\(\exists x(x\ne y)\) is true for every \(y\), whereas
\(c\ne y\) is false when \(y=c\).
\end{proof}
\end{annexitem}

\begin{annexitem}{annex:p12.opt-optimized-skolem}{thm:p12-opt-optimized-skolem}{The two witness cases}
\begin{proof}
Let \(M\models F\). For each tuple \(\bar d\) of values for \(\bar y\),
if \(\exists\bar x(A\land B)\) is true choose a simultaneous
witness tuple for \(A\land B\). If it is false choose a witness
for \(A\), which exists by the premise
\(F\models\forall\bar y\exists\bar x A\).
Use its \(i\)-th coordinate as \(f_i^{M'}(\bar d)\).
This defines total functions; all formulas are evaluated in \(M\)
before interpreting the fresh symbols. The choices ensure
\(\forall\bar y A\sigma\), and also ensure at every assignment
that \(E\) implies \(B\sigma\). Since its occurrence is positive
in NNF, contextual monotonicity gives \(F[\pi/B\sigma]\).
This is the asserted expansion.

Conversely take a model \(M'\) of the transformed sentence.
At every tuple \(\bar d\), \(A\sigma\) holds. If \(B\sigma\)
also holds, the tuple \((f_i^{M'}(\bar d))_{i=1}^k\) witnesses
\(E\) in the reduct. Hence \(B\sigma\to E\) at all assignments,
and contextual monotonicity proves \(F\).
This direction uses no choice and no unproved inference about
the truth of \(B\) where \(E\) was false.
\end{proof}
\end{annexitem}

\begin{annexitem}{annex:p12.opt-selector}{thm:p12-opt-selector}{Empty and nonempty fibers}
\begin{proof}
For each parameter tuple \(\bar b\), define
\[
 S(\bar b)=\{\bar a\in D^k:M[\bar y/\bar b,\bar x/\bar a]\models A\}.
\]
Choose \(d_0\in D\). At a nonempty fiber choose
\(\bar a_0(\bar b)\in S(\bar b)\), and set
\[
 h(\bar b,\bar w)=
 \begin{cases}
 \bar w,&\bar w\in S(\bar b),\\
 \bar a_0(\bar b),&\bar w\notin S(\bar b),\ S(\bar b)\ne\varnothing,\\
 (d_0,\ldots,d_0),&S(\bar b)=\varnothing.
 \end{cases}
\]
Its coordinate projections interpret the fresh \(h_j\).
On a nonempty fiber this is a retraction, so all its values satisfy
\(A\) and it fixes every tuple that satisfies \(A\).
Given a witness \(\bar a\) for \(A\land B\), take
\(\bar v=\bar a\); then \(h(\bar b,\bar v)=\bar a\), preserving
\(B\), and every \(\bar w\) gives \(A\).
Other free variables of \(B\) are simply left at their original
values. Thus the implication is uniform over all assignments,
although its antecedent is false at an empty fiber.
\end{proof}
\end{annexitem}

\begin{annexitem}{annex:p12.opt-strong-skolem}{thm:p12-opt-strong-skolem}{Two complete uniform constructions}
\begin{proof}
First note that \(\operatorname{FV}(A_i)\setminus\{\bar x\}\)
is contained in \(Z_1\cup\cdots\cup Z_i\).
The proof works more generally for disjoint ordered parameter blocks
with this containment property, even if a block contains parameters
unused by that conjunct. We prove this slightly more general uniform
statement to justify the source's intermediate padded parameter lists.

\emph{Selector-composition construction.}
Induct on \(n\). For \(n=1\), at every parameter tuple with a
witness choose a witness; elsewhere choose \((d_0,\ldots,d_0)\).
The coordinates define \(f_j(Z_1)\), establishing the uniform
implication. Suppose \(n>1\). Apply the selector construction
to \(A_1\), with parameters \(Z_1\), obtaining the vector function
\(h(Z_1,\bar v)\). Put
\[
 B_i=A_i\{\bar x\mapsto h(Z_1,\bar v)\}\quad(2\le i\le n).
\]
Use parameter blocks \(Z_1Z_2,Z_3,\ldots,Z_n\) for these \(n-1\)
conjuncts. The containment condition holds: substitution only adds
the variables in \(Z_1\) and the existential tuple \(\bar v\).
This remains true when an \(A_i\) contains no variable from
\(\bar x\), which is why exact equality of lexical free-variable
sets is unnecessary.

In the already expanded structure, the induction hypothesis supplies
total functions \(g_1,\ldots,g_k\) for the residual existential
\(\exists\bar v(B_2\land\cdots\land B_n)\), uniformly over
all full parameter tuples. Define, for every tuple \(z\),
\[
 f_j(z)=h_j(z_1,g_1(z),\ldots,g_k(z)).
\]
Now fix a tuple \(z\) at which \(E\) holds. Its witness lies in
the nonempty \(A_1\)-fiber. Retraction fixes that witness, so the
residual existential holds at \(z\).
The induction hypothesis yields, for every \(i\ge2\) and every
replacement of its later parameter blocks by \(w_{i+1},\ldots,w_n\),
the truth of \(B_i\) with \(\bar v\) interpreted by
\(g(z_1,\ldots,z_i,w_{i+1},\ldots,w_n)\).
After the definition of \(f\), this is precisely conjunct \(i\)
of \(\Xi\). For conjunct \(1\), its nonempty fiber makes
\(A_1(h(z_1,\bar v))\) true for \emph{every} \(\bar v\).
In particular it is true for every value of
\(g(z_1,w_2,\ldots,w_n)\). This proves the first conjunct of
\(\Xi\). The construction was made once for the entire structure,
not separately for assignments later selected by an enclosing
quantifier. Finally discard the auxiliary \(h,g\) symbols, retaining
the composed \(f\)'s. This completes the induction and the source's
composition route with totality at empty fibers made explicit.

\emph{Direct prefix-witness construction.}
For a full parameter tuple \(z=(z_1,\ldots,z_n)\), let
\[
 S_i(z_1,\ldots,z_i)=
 \{\bar a:M[\bar x/\bar a,Z_{\le i}/z_{\le i}]
                     \models A_1\land\cdots\land A_i\},
\]
and let \(S_0=D^k\). Evaluation is well defined because the first
\(i\) conjuncts use only the displayed parameters. For fixed \(z\),
these sets are nested decreasing as \(i\) increases.
Let \(j(z)\) be the largest \(i\in\{0,\ldots,n\}\) with
\(S_i(z_{\le i})\ne\varnothing\).
Choose \(f(z)\in S_{j(z)}(z_{\le j(z)})\).
These choices define total functions coordinatewise.

If \(E\) is true at \(z\), each \(S_i(z_{\le i})\) is nonempty.
Fix \(i\) and replace the later blocks of \(z\) arbitrarily to
obtain \(z'\). The \(i\)-th prefix set is unchanged and nonempty,
so \(j(z')\ge i\). The chosen \(f(z')\) lies in
\(S_{j(z')}(z'_{\le j(z')})\subseteq S_i(z_{\le i})\).
It therefore satisfies \(A_i\), which is the universal assertion
in conjunct \(i\) of \(\Xi\). This holds for every \(i\) and
every later tuple, proving the same uniform implication directly.

\emph{Reduct and surrounding context.}
In any expansion satisfying \(\Xi\) at \(z\), instantiate all its
later universal blocks by the corresponding actual blocks of \(z\).
Each conjunct is then true for the single tuple \(f(z)\);
that tuple witnesses \(E\). This direction does not use either
witness construction. At an arbitrary positive occurrence in an NNF
sentence, the uniform forward implication preserves truth in the
chosen expansion; the uniform reverse implication preserves truth
in the reduct. Contextual monotonicity includes every surrounding
universal and existential binder, so there is no restriction to
an outermost existential or to finite domains.
\end{proof}
\end{annexitem}

\begin{annexitem}{annex:p12.opt-equality}{thm:p12-opt-equality}{Equality, capture avoidance and cardinality}
\begin{proof}
Fix \(M,s\), and let \(a=t^{M,s}\).
Since \(x\) does not occur in \(t\), its value is unchanged when
the assignment to \(x\) changes. Thus
\(M,s[x/b]\models x=t\) exactly when \(b=a\).
The existential formula therefore holds exactly when
\(M,s[x/a]\models A\), and the universal implication has the
same condition: all other \(b\)'s have a false antecedent.
Capture-free substitution gives
\[
 M,s[x/a]\models A\quad\Longleftrightarrow\quad M,s\models A[x/t].
\]
For completeness, this substitution fact follows from term induction
for atoms and from preservation of truth by each Boolean operation.
At a binder \(\forall z\) or \(\exists z\), first choose \(z\)
different from \(x\) and all variables of \(t\); assignment to
\(z\) then commutes with changing \(x\) to \(a\).
Apply the induction hypothesis for each value of \(z\), and use
respectively all values or one witnessing value. A binder for \(x\)
itself protects its occurrences and is unchanged. This proves the
substitution fact in every case.

Identity makes \(t=t\) true. If \(y\) is absent from \(t\), put
\(a=t^{M,s}\); choosing \(y=a\) refutes \(\forall y(t\ne y)\)
and witnesses \(\exists y(t=y)\).
On a singleton domain every \(y\) equals \(a\); on a larger domain
choose \(b\ne a\), which refutes the universal equality and
witnesses the existential inequality. These prove all six
cardinality cases. Contextual equivalence gives replacement on
any class on which the same conditions hold on both sides.

For termination, equality substitution removes its quantified
binder and inserts only terms, hence reduces the total number
of quantifiers. Reflexive and cardinality reductions cannot
increase that number and decrease size if they preserve it.
The standard Boolean simplifications do not increase quantifier
count and decrease size when it is unchanged.
The lexicographically ordered pair (quantifier count, formula size)
therefore decreases at every such step and is well founded on
nonnegative integers. Term size alone would not justify termination
of the substitution step.
\end{proof}
\end{annexitem}

\begin{annexitem}{annex:p12.opt-atom-expansion}{thm:p12-opt-atom-expansion}{Retained definitions and guards}
\begin{proof}
In a model of \(\Theta\), universal instantiation of its implied
definition at \(\sigma\) gives equality of the truth values of
\(P(\bar s)\) and \(A\sigma\) at every assignment.
The substitution induction in the preceding proofs licenses that
instantiation. Contextual equivalence proves replacement for any
polarity. A globally established guard gives the same equivalence.
For the guarded-context interface, fix the outside assignment.
If the guard is false, the guarded implication is true on both
sides. If it is true, the conditional definition supplies the
needed equivalence throughout the occurrence's context, provided
no intervening binder changes a free variable of the guard.
This last scope condition is required when the guard is outside
that binder; alternatively establish its universal closure as in
the first interface. Applying contextual equivalence then proves
the assertion in the true-guard case.

With one \(P\)-free definition body, substitution inserts no
occurrence of \(P\). Each expansion deletes one targeted occurrence,
so their finite count strictly decreases.
For elimination, interpret \(P(\bar d)\) in a reduct by the truth
of \(A\) at \(\bar x=\bar d\). The distinct head variables let
every tuple have that interpretation, and \(P\)-freeness means the
truth of \(A\) was already defined in the reduct. The resulting
expansion satisfies the defining equivalence, hence every other
formula exactly when its expanded version holds.
Conversely the defining equivalence in any original model forces
this same interpretation of \(P\) for every tuple. This proves
existence and uniqueness of the expansion and the claimed set of
reducts. Retaining a recursively self-referential definition does
not supply this uniqueness argument.
\end{proof}
\end{annexitem}

\begin{annexitem}{annex:p12.opt-bdd}{thm:p12-opt-bdd}{Splitting, canonical reduction and clauses}
\begin{proof}
Fix a valuation. If \(A\) is true, replacement of its designated
occurrence by \(\top\) leaves the value of \(F\) unchanged;
the first implication on the right therefore has value \(F\),
and the second has false antecedent and is true.
If \(A\) is false the roles reverse, with replacement by \(\bot\).
These two exhaustive cases prove the splitting identity even
when other occurrences of \(A\) remain in \(F\).

For complete splitting choose \(v_1\), substitute its value at
all its occurrences, and recurse on the remaining ordered variables.
At each recursion the number of available variables decreases,
ending in a constant after at most \(m\) decisions.
Its full tree has at most \(1+2+\cdots+2^m=2^{m+1}-1\)
nodes. The truth computed by the tree follows the substituted
valuation and equals the formula's truth.
Merging identical nodes changes no successor or decision value.
If both successors agree, testing that variable changes no value,
so deletion is also sound. Each reduction decreases the number
of vertices, making repeated reduction terminate.
Discarding unreachable vertices preserves every path from the root
and hence preserves the represented function. It also ensures that
the graph being compared contains no unrelated components.

For uniqueness induct on the number of available variables.
A constant function has only its terminal in a reduced graph:
by the induction applied to each root child, both children of
any node representing a constant would be that terminal and the
node would have been removed.
For a nonconstant function \(f\), its root variable must be the
least variable on which \(f\) actually depends. Indeed a later
root would skip a dependency; an earlier root with no dependency
would have equal restricted functions on its two branches.
By induction their reduced child graphs would be identical,
forcing deletion of that root. At the required least variable,
the low and high functions are the uniquely determined restrictions
of \(f\), and induction determines both child graphs. Sharing
identical triples is forced by reduction. Thus every node and
edge are determined up to the stated isomorphism.

A valuation reaches \(0\) exactly when its branch choices extend
one of the root-to-\(0\) paths. The clause associated with that
path is false exactly for those extensions: each of its literals
negates the branch choice. Hence all path clauses hold exactly
when no such path is followed, that is, when the diagram evaluates
to \(1\). A \(0\)-terminal at the root contributes the empty
clause; a \(1\)-terminal with no false path contributes the empty
conjunction. These cover both degenerate functions.
\end{proof}
\end{annexitem}
